% GENERATED FILE. Edit canonical lessons, then run npm run generate:books.
% canonical-source-hash: fnv1a64:878a6cb5949c0c20
% canonical-lessons: ML-C128-house, ML-C128-room, ML-C128-kitchen, ML-C128-bathroom, ML-C128-yard

\chapter{Come In, This Is Home}
\label{ch:ml-come-in-this-is-home}
\hlchaptermodality{128}

\begin{chapteropening}
\textbf{By the end of this chapter:} \emph{I can name a house, a room, the kitchen, the bathroom and the yard.}

Complete the last lesson of chapter 128: I can name a house, a room, the kitchen, the bathroom and the yard.
\end{chapteropening}

\section[\texorpdfstring{\ml{വീട്} (vīṭŭ)}{വീട് (vīṭŭ)}]{\ml{വീട്} (vīṭŭ) — a house, home}

\label{lesson:ML-C128-house}



\begin{quote}
Before the new one: say the Malayalam for a holiday, then the Malayalam for I'll call.
\end{quote}

\subsection*{You'll want to know: \ml{വീട്}}
\textbf{\ml{വീട്}} (\emph{vīṭŭ}) — \textquotedblleft{}a house\textquotedblright{}, and also \textquotedblleft{}home\textquotedblright{}.

\textbf{\ml{എന്റെ} \ml{വീട്}} (\emph{enṟe vīṭŭ}) — \textquotedblleft{}my house\textquotedblright{}, with \textbf{\ml{എന്റെ}}, \textquotedblleft{}my\textquotedblright{}, which you already know.

\begin{cousinweb}
\textbf{\ml{വീട്}} (\emph{vīṭŭ}) is native Dravidian: DEDR 5393 lists it, \textquotedblleft{}house\textquotedblright{}, under the verb \textquotedblleft{}to let go, leave off\textquotedblright{}, with Tamil \emph{viṭu}. DEDR 5393 notes Sanskrit \emph{avamocana}, \textquotedblleft{}inn\textquotedblright{}, as a translation of Tamil \emph{viṭuti}.
\end{cousinweb}

\begin{grammarlens}[title={What you've built}]
The place every visit starts.
\end{grammarlens}

\subsection*{Your turn}
\begin{itemize}
\raggedright
  \item \emph{Say it:} \emph{vīṭŭ}
  \item \emph{Say it:} \emph{vīṭŭ}, once more
  \item \emph{Say it:} say \emph{enṟe vīṭŭ}
  \item \emph{Recall:} say the Malayalam for a holiday, then the Malayalam for I'll call, then say \emph{vīṭŭ} again
\end{itemize}

\begin{tcolorbox}[breakable,colback=teal!4,colframe=teal!35!black,title={Before you move on}]
What does \ml{വീട്} mean? (\textquotedblleft{}A house, home\textquotedblright{}.) Say it once more.
\end{tcolorbox}
\section[\texorpdfstring{\ml{മുറി} (muṟi)}{മുറി (muṟi)}]{\ml{മുറി} (muṟi) — a room}

\label{lesson:ML-C128-room}



\begin{quote}
Before the new one: say the Malayalam for I'll call, then the Malayalam for a house.
\end{quote}

\subsection*{You'll want to know: \ml{മുറി}}
\textbf{\ml{മുറി}} (\emph{muṟi}) — \textquotedblleft{}a room\textquotedblright{}. \textbf{\ml{ഇത്} \ml{എന്റെ} \ml{മുറി}} — \textquotedblleft{}this is my room\textquotedblright{}.

\begin{cousinweb}
\textbf{\ml{മുറി}} (\emph{muṟi}) is native Dravidian, one word with Tamil \emph{muṟi} and Kannada \emph{muṟi} (DEDR 5008). Its first sense is \textquotedblleft{}a piece, a fragment\textquotedblright{}: a room is a part cut off from a house.
\end{cousinweb}

\begin{grammarlens}[title={What you've built}]
One part of the house.
\end{grammarlens}

\subsection*{Your turn}
\begin{itemize}
\raggedright
  \item \emph{Say it:} \emph{muṟi}
  \item \emph{Say it:} \emph{muṟi}, once more
  \item \emph{Say it:} say \emph{itŭ enṟe muṟi}
  \item \emph{Recall:} say the Malayalam for I'll call, then the Malayalam for a house, then say \emph{muṟi} again
\end{itemize}

\begin{tcolorbox}[breakable,colback=teal!4,colframe=teal!35!black,title={Before you move on}]
What does \ml{മുറി} mean? (\textquotedblleft{}A room\textquotedblright{}.) Say it once more.
\end{tcolorbox}
\section[\texorpdfstring{\ml{അടുക്കള} (aṭukkaḷa)}{അടുക്കള (aṭukkaḷa)}]{\ml{അടുക്കള} (aṭukkaḷa) — the kitchen}

\label{lesson:ML-C128-kitchen}



\begin{quote}
Before the new one: say the Malayalam for a house, then the Malayalam for a room.
\end{quote}

\subsection*{You'll want to know: \ml{അടുക്കള}}
\textbf{\ml{അടുക്കള}} (\emph{aṭukkaḷa}) — \textquotedblleft{}the kitchen\textquotedblright{}, where the \textbf{\ml{ചോറ്}} is cooked.

\begin{cousinweb}
\textbf{\ml{അടുക്കള}} (\emph{aṭukkaḷa}) is native Dravidian, built on the verb \textquotedblleft{}to cook\textquotedblright{}; DEDR 76 lists Malayalam \emph{aṭukkaḷa} itself, beside Tamil \emph{aṭu} and Kannada \emph{aḍu}. DEDR 76 compares Sanskrit \emph{aṭṭa}, \textquotedblleft{}dried\textquotedblright{}.
\end{cousinweb}

\begin{grammarlens}[title={What you've built}]
The room with the fire.
\end{grammarlens}

\subsection*{Your turn}
\begin{itemize}
\raggedright
  \item \emph{Say it:} \emph{aṭukkaḷa}
  \item \emph{Say it:} \emph{aṭukkaḷa}, once more
  \item \emph{Say it:} ask \emph{aṭukkaḷa evide?}
  \item \emph{Recall:} say the Malayalam for a house, then the Malayalam for a room, then say \emph{aṭukkaḷa} again
\end{itemize}

\begin{tcolorbox}[breakable,colback=teal!4,colframe=teal!35!black,title={Before you move on}]
What does \ml{അടുക്കള} mean? (\textquotedblleft{}The kitchen\textquotedblright{}.) Say it once more.
\end{tcolorbox}
\section[\texorpdfstring{\ml{കുളിമുറി} (kuḷimuṟi)}{കുളിമുറി (kuḷimuṟi)}]{\ml{കുളിമുറി} (kuḷimuṟi) — the bathroom}

\label{lesson:ML-C128-bathroom}



\begin{quote}
Before the new one: say the Malayalam for a room, then the Malayalam for the kitchen.
\end{quote}

\subsection*{You'll want to know: \ml{കുളിമുറി}}
\textbf{\ml{കുളിമുറി}} (\emph{kuḷimuṟi}) — \textquotedblleft{}the bathroom\textquotedblright{}: \textbf{\ml{കുളി}} (\emph{kuḷi}), \textquotedblleft{}a bath\textquotedblright{}, with \textbf{\ml{മുറി}}, \textquotedblleft{}room\textquotedblright{}.

\begin{cousinweb}
\textbf{\ml{കുളിമുറി}} (\emph{kuḷimuṟi}) is native Dravidian: \emph{kuḷi}, \textquotedblleft{}bathing\textquotedblright{} (DEDR 1832, with Tamil \emph{kuḷi}), with \emph{muṟi}, \textquotedblleft{}room\textquotedblright{} (DEDR 5008).
\end{cousinweb}

\begin{grammarlens}[title={What you've built}]
A room word you can take apart.
\end{grammarlens}

\subsection*{Your turn}
\begin{itemize}
\raggedright
  \item \emph{Say it:} \emph{kuḷimuṟi}
  \item \emph{Say it:} \emph{kuḷimuṟi}, once more
  \item \emph{Say it:} ask \emph{kuḷimuṟi evide?}
  \item \emph{Recall:} say the Malayalam for a room, then the Malayalam for the kitchen, then say \emph{kuḷimuṟi} again
\end{itemize}

\begin{tcolorbox}[breakable,colback=teal!4,colframe=teal!35!black,title={Before you move on}]
What does \ml{കുളിമുറി} mean? (\textquotedblleft{}The bathroom\textquotedblright{}.) Say it once more.
\end{tcolorbox}
\section[\texorpdfstring{\ml{മുറ്റം} (muṟṟaṁ)}{മുറ്റം (muṟṟaṁ)}]{\ml{മുറ്റം} (muṟṟaṁ) — the yard}

\label{lesson:ML-C128-yard}



\begin{quote}
Before the new one: say the Malayalam for the kitchen, then the Malayalam for the bathroom.
\end{quote}

\subsection*{You'll want to know: \ml{മുറ്റം}}
\textbf{\ml{മുറ്റം}} (\emph{muṟṟaṁ}) — \textquotedblleft{}the yard\textquotedblright{}, the swept ground in front of a Kerala house where the \textbf{\ml{കോലം}} is drawn.

\begin{cousinweb}
\textbf{\ml{മുറ്റം}} (\emph{muṟṟaṁ}) is native Dravidian, one word with Tamil \emph{muṟṟam} (DEDR 5016).
\end{cousinweb}

\begin{grammarlens}[title={What you've built}]
Outside, but still home.
\end{grammarlens}

\subsection*{Your turn}
\begin{itemize}
\raggedright
  \item \emph{Say it:} \emph{muṟṟaṁ}
  \item \emph{Say it:} \emph{muṟṟaṁ}, once more
  \item \emph{Say it:} say \emph{muṟṟaṁ}
  \item \emph{Recall:} say the Malayalam for the kitchen, then the Malayalam for the bathroom, then say \emph{muṟṟaṁ} again
\end{itemize}

\begin{tcolorbox}[breakable,colback=teal!4,colframe=teal!35!black,title={Before you move on}]
What does \ml{മുറ്റം} mean? (\textquotedblleft{}The yard\textquotedblright{}.) Say it once more.
\end{tcolorbox}
